\chapter{Conclusions and Future Work}
\label{s_conclusions}

\section{Conclusions}
This thesis has presented the design, implementation, and evaluation of a
microservice-based path planning framework for urban air mobility (UAM)
operations. By decomposing the system into independent services, the planner
was able to integrate heterogeneous components implemented in different
languages and tools—namely, the A* global planner (Python), the CFMS trajectory
generator (C++), and the NDMap conflict detection module (C++). This modular
design demonstrates the potential of microservices for building scalable and
extensible architectures for U-space operations.

The framework was evaluated using a representative vertiport and corridor
network modeled over the metropolitan area of Hamburg. Experiments focused on
quantifying delays at different processing stages, including A* path finding,
trajectory generation, and overall end-to-end response time. Results showed
that:
\begin{itemize}
    \item The A* path finding step contributed negligibly to the total delay,
    with runtimes consistently below 1~ms.
    \item Trajectory generation dominated the computational cost, with runtimes
    scaling linearly with path distance.
    \item The total processing time per trajectory request was on the order of
    200~ms, demonstrating the feasibility of near real-time operation for
    single-request scenarios.
    \item Conflict detection and resolution were successfully demonstrated using
    NDMap, with lateral rerouting as the chosen strategy to avoid overlapping
    trajectories.
\end{itemize}

Overall, the system provides a primitive but functional prototype of a path
planner for UAM operations, highlighting both the feasibility and the
challenges of integrating multiple specialized modules into a unified framework.

\section{Future Work}
While the present study establishes a foundation, several extensions are
necessary to bring the framework closer to realistic deployment:

\subsection*{Enhanced Conflict Resolution}
The current work focused exclusively on lateral rerouting through penalization
of conflicting edges in the corridor graph. Future extensions should include
additional resolution strategies such as:
\begin{itemize}
    \item \textbf{Temporal deconfliction:} delaying departure times to avoid
    simultaneous use of a corridor.
    \item \textbf{Vertical deconfliction:} assigning different altitude layers
    to conflicting trajectories.
    \item \textbf{Hybrid strategies:} combining lateral, temporal, and vertical
    maneuvers to balance efficiency and safety.
\end{itemize}

\subsection*{Complex Airspace Modeling}
The current corridor network was static, uniform in altitude, and of fixed
width. Future studies should incorporate:
\begin{itemize}
    \item Multiple altitude layers with dynamic assignment,
    \item Variable corridor widths depending on traffic density or risk,
    \item More realistic urban scenarios with larger networks and higher traffic
    volumes.
\end{itemize}

\subsection*{Incorporation of External Factors}
To better approximate real-world operations, future versions should integrate:
\begin{itemize}
    \item \textbf{Weather effects:} wind, turbulence, and no-fly zones affecting
    edge weights and feasibility.
    \item \textbf{Risk maps:} penalizing corridors near sensitive areas (schools,
    hospitals, etc.) to reflect societal and regulatory constraints.
    \item \textbf{Dynamic re-weighting:} adjusting corridor costs in real time
    based on environmental and operational factors.
\end{itemize}

\subsection*{Aircraft Performance Models}
The present study employed simplified kinematic constraints. For improved
fidelity, future work should model:
\begin{itemize}
    \item Performance envelopes specific to different eVTOL types,
    \item More detailed climb, descent, and turn dynamics,
    \item Energy consumption models for battery-electric aircraft.
\end{itemize}

\subsection*{Scalability and Multi-Agent Evaluation}
The synchronous evaluation with 100 independent trajectory requests provided
baseline performance characteristics. Future studies should investigate:
\begin{itemize}
    \item Asynchronous and parallel trajectory requests at scale,
    \item End-to-end system latency under high traffic densities.
\end{itemize}

\medskip
\noindent
In summary, this thesis demonstrates that a microservice-based approach to UAM
path planning is both technically feasible and computationally efficient for
baseline scenarios. Expanding the framework with more advanced conflict
resolution strategies, richer airspace and aircraft models, and larger-scale
evaluations will be critical steps toward enabling safe, scalable, and
realistic integration of UAM operations into the future air traffic management
ecosystem.
